\documentclass[11pt]{article}
\usepackage[T1]{fontenc}
\usepackage[utf8]{inputenc}
\usepackage{amsmath,amssymb,amsthm,mathtools,hyperref}
\usepackage[margin=25mm]{geometry}
\newcommand{\Hom}{\operatorname{Hom}}
\newcommand{\Spec}{\operatorname{Spec}}
\newcommand{\Tor}{\operatorname{Tor}}
\newcommand{\Ext}{\operatorname{Ext}}
\newcommand{\End}{\operatorname{End}}
\newcommand{\im}{\operatorname{Im}}
\newcommand{\coker}{\operatorname{Coker}}
\newcommand{\rk}{\operatorname{rk}}
\newcommand{\Supp}{\operatorname{Supp}}
\newcommand{\colim}{\mathop{\rm colim}}
\newcommand{\m}{\mathfrak m}
\newcommand{\p}{\mathfrak p}
\newcommand{\q}{\mathfrak q}
\title{WTB\_94\_002--013: proofs, corrections, and certified background}
\author{GPT-6 Astra Pro; single mathematical stream}
\date{19 November 2026}
\begin{document}
\maketitle

\section*{Conventions and separation of evidence}
Manuscript locators use scan numbers, followed where useful by lines of witness C.
The scan is authoritative. A quotation in this file is only a short anchor;
the complete modern statement is explicitly a formalization.
Rings are commutative with identity, modules unital. \emph{Exact à gauche}
for the initially set-valued functor means preservation of finite limits.
Such a functor has canonical additive and $A$-linear structures: transport
the zero, addition, inverse, and scalar-operation diagrams through its
preserved finite products. Naturality supplies these structures on maps.
Finite means finitely generated unless length is expressly specified.
Direct limits in these criteria mean filtered colimits; this is explicit on
scans 3, 4 and 6 and reconstructed on scan 5, where only the arrow is written.

Statements identified as SETTLED are certified input, carried without a new
proof as the prompt requires. A prior run's discarded status is not the
settled verdict. A section that cites this input does not claim that the
current run independently proved it. New proofs are labeled ours; separate
verification records identify the actual second checks.

\section{94.C01: change of rings}
\label{C01}
\textbf{His anchor:} scan 3, C:8--15, \emph{$F_B$ l'est par $P_B$}.
\textbf{Modern statement (ours):} if $F(M)=\Hom_A(P,M)$ and $A\to B$
is a ring map, its restriction to $B$-modules is represented by $B\otimes_A P$.
No flatness or finiteness is needed.

\textbf{Proof (ours).}
Send a $B$-linear map $g:B\otimes_A P\to M$ to $p\mapsto g(1\otimes p)$.
The inverse sends $u:P\to M$ to $b\otimes p\mapsto bu(p)$.
The balancing identity proves well-definedness, and the two composites
are identities. Both maps are natural in $M$. Localization is the
special case $B=A_f$. This also proves compatibility under successive
base changes, which will provide the descent cocycle below.

\section{94.C02: Lemme 1, faithful-flat descent}
\label{C02}
\textbf{His anchor:} scan 3, C:17--24, \emph{$F$ exact à gauche}.
\textbf{Modern statement (ours):} for faithfully flat $A\to B$,
$F:\mathrm{Mod}(A)\to\mathrm{Set}$ is representable if and only if it is
left exact and $F_B$ is representable. The equivalence remains true with
represented by a finite module, or by a finitely presented module,
on both sides.

\textbf{Proof (ours, using the verified descent theorem).}
The forward implication follows from C01 and preservation of limits by
$\Hom(P,-)$. Conversely let $Q$ represent $F_B$. Put $C=B\otimes_A B$.
The two base changes of $Q$ to $C$ both represent exactly the restriction
of $F$ to $C$-modules. Yoneda gives a unique isomorphism compatible with
these identifications. Over $B\otimes_A B\otimes_A B$ the two possible
composites induce the same identity transformation of $F$; hence they
are equal. This is a descent datum, effective by
Stacks Proposition 35.3.9, Tag 023N \cite{Stacks023N}. Thus
$Q\simeq B\otimes_A P$ compatibly with descent for an $A$-module $P$.

For every $M$, faithful flatness gives the equalizer
\[
 M\longrightarrow B\otimes_A M
 \rightrightarrows B\otimes_A B\otimes_A M .
\]
Applying $F$ preserves it. Under the representing identifications the
two arrows become the descent arrows for $\Hom_A(P,M)$, so their
equalizer is $\Hom_A(P,M)$. This constructs the required natural
isomorphism, rather than just an object with the right scalar extension.

For completeness, finiteness descends as follows. A finite set of
generators of $B\otimes P$ uses only finitely many elements of $P$.
Their span $P_0$ has $B\otimes(P/P_0)=0$ and hence $P=P_0$.
If $B\otimes P$ is finitely presented, choose $A^r\twoheadrightarrow P$
with kernel $K$. Flatness identifies $B\otimes K$ with the kernel of
$B^r\to B\otimes P$, which is finite; the preceding argument makes $K$
finite. The converse ascents follow from tensoring a presentation.
These two descent facts also match Stacks Lemma 10.83.2, Tag 03C4
\cite{Stacks03C4}.

\section{94.C03: Lemme 2, principal completion and localization}
\label{C03}
\textbf{His anchor:} scans 3--4, \emph{$M$ un $\widehat A$-module}.
\textbf{Modern statement (ours):} let $A$ be noetherian, $f\in A$,
$R=\varprojlim_n A/f^n A$. Then $F$ is represented by a finite $A$-module
if and only if:
(i) $F$ is left exact;
(ii) for every finite $R$-module $M$ the canonical map
$F(M)\to\varprojlim_n F(M/f^nM)$ is bijective;
(iii) $F$ commutes with filtered colimits;
(iv) $F_{A_f}$ is representable; and
(v) each $F_{A/f^nA}$ is representable, $n\geq1$.
The hat on the argument of $F$ in the earlier wording of (ii) is canceled
on scan 3; the completed coefficient ring is not canceled.

\textbf{Necessity (ours).}
A finite module over a noetherian ring is finitely presented.
Its Hom functor preserves filtered colimits, as is seen by applying
$\ker(M^r\to M^s)$ to a finite presentation and using the commutation
of filtered colimits with finite limits. It preserves all limits, and
a finite $R$-module is $f$-adically complete. C01 supplies (iv)--(v).

\textbf{Sufficiency (ours).}
Choose a representing finite module $Q_n$ for $F_{A/f^nA}$.
It is finite because its Hom functor commutes with filtered colimits:
write $Q_n$ as the filtered union of its finite submodules, and factor
its identity through one of them. Yoneda and C01 give compatible
isomorphisms $Q_{n+1}/f^n Q_{n+1}\simeq Q_n$.
The finite formal-module equivalence, Stacks Lemma 30.23.1, Tag 087W
\cite{Stacks087W}, yields a finite $R$-module $Q=\varprojlim Q_n$
with $Q/f^nQ\simeq Q_n$. Explicitly, lift generators of $Q_1$
compatibly; they generate every $Q_n$ by Nakayama. The kernels of the
resulting maps $(A/f^n)^r\to Q_n$ have surjective transitions. Taking
limits gives a quotient of $R^r$ with precisely these reductions.

For finite $R$-modules $M$ we consequently have, naturally,
\[
 F(M)\simeq\varprojlim_n\Hom_{A/f^n}(Q_n,M/f^nM)
 \simeq\Hom_R(Q,M).
\]
The second bijection follows by taking compatible maps into the
complete module $M$; $R$-linearity follows at every finite quotient.
Since $R$ is noetherian, $Q$ is finitely presented. Both sides preserve
filtered colimits, so the isomorphism extends to all $R$-modules.

The ring map $A\to R\times A_f$ is faithfully flat. Flatness of $A\to R$
is Stacks Lemma 10.97.2 \cite{StacksCompletion}. Every prime containing
$f$ lifts through $R/fR=A/fA$; primes not containing $f$ are covered by
$\Spec A_f$. Thus the flat map is surjective on spectra.
Its restricted functor is represented by the pair consisting of $Q$
and the representative on $A_f$; preservation of products is used here.
C02 descends it. Finally (iii), applied to the filtered union of finite
submodules of the descended representative, makes that representative
finite. This also handles nilpotent or invertible $f$ and the zero ring.

\section{94.C03b: a local wording slip in the proof}
\label{C03b}
Scan 4, C:4 writes \emph{$A$-modules de t.f.} without a visible hat.
The preceding construction only proves representation on finite
$\widehat A$-modules; these are exactly the objects used in (ii).
The correct bridge to all $\widehat A$-modules is their filtered union
of finite $\widehat A$-submodules. The two categories cannot be
identified: for $A=\mathbb Z$ and $f=p$, $\mathbb Z_p$ is cyclic over
itself but not finite over $\mathbb Z$ (a finite $\mathbb Z$-module is
countable; $\mathbb Z_p$ has distinct expansions for all infinite
$p$-adic digit strings). This is an authorial shorthand/slip, not a
transcription correction. The proof just given supplies the bridge.

\section{94.C04: Proposition 3 --- certified background}
\label{C04}
\textbf{His anchor:} scans 4--5, \emph{quotient non nul}.
With $A$ noetherian, finite representability is equivalent to
left exactness, commutation with filtered colimits, the completion
comparison of C03 for every $f\in A$, and the following condition:
for every nonzero quotient $B$ of $A$, some nonnilpotent $g\in B$
makes $F_{B_g}$ representable.

This is SETTLED statement H035.5, carried as proved background, not
reproved here. The prior label \emph{false as written} was explicitly
refuted by the supplied check: the nonzero qualifier is in the ink.
Our new image check confirms the leader attaching it to $B$.
The written proof uses noetherian induction on quotient ideals and
C03; the step $g^{n+1}\ne0$ is visible on scan 5.
There is no missing nonzero hypothesis to add to this proposition.

\section{94.C05: Proposition 4 and its Artinian step}
\label{C05}
\textbf{His anchor:} scans 5--6, \emph{l'injectif indécomposable $I_n$}.
\textbf{Modern statement (ours):} for a noetherian local $(A,\m)$,
$F$ is represented by a finite $A$-module iff it is left exact,
preserves filtered colimits, satisfies
\[
 F(M)\simeq\varprojlim_n F(M/\m^{n+1}M)
\]
for each finite $\widehat A$-module $M$, and takes such modules to finite
$\widehat A$-modules. It suffices to impose the last condition on
modules finite over $A$, as his parenthesis says. In fact finite length
modules, and even the single module $A/\m$, suffice (RD01 below).
The live $\widehat M$ in the formula is harmless: $M$ is already complete.

\textbf{Artinian step (ours).}
Let $(R,\m,k)$ be Artinian local, $E=E_R(k)$, and
$D=\Hom_R(-,E)$. The verified finite-length duality gives finite $E$,
$D^2\simeq1$, and $\operatorname{End}_R(E)=R$
(Stacks Lemmas 47.6.1--3 \cite{StacksArtinDuality}).
Suppose $G$ is $R$-linear and left exact on finite $R$-modules, with
the canonical scalar structure described above, and $G(E)$ is
finite. Set $P=D(G(E))$. There is a natural map
\[
 G(M)\longrightarrow\Hom_R(DM,G(E)),\qquad
 x\longmapsto\bigl(\lambda\mapsto G(\lambda)(x)\bigr).
\]
It is an isomorphism for $M=E^r$, since $\operatorname{End}(E)=R$.
Every finite $M$ has an exact copresentation
$0\to M\to E^r\to E^s$: dualize a finite free presentation of $DM$.
Apply both left-exact functors to this copresentation. The two
isomorphisms at $E^r,E^s$ identify their kernels, proving the displayed
map is an isomorphism for $M$. Duality changes its target to
$\Hom_R(P,M)$. This proves representability, including naturality.

\textbf{Local proof (ours).}
Restriction to $\widehat A$ preserves all the stated assumptions
needed below: finite $A/\m^{n+1}$-modules are finite over $A$.
The completion is faithfully flat (Stacks Lemma 10.97.3).
For $A_n=A/\m^{n+1}$ the Artinian step gives a finite representative
$P_n$; (iv) applies to its finite injective cogenerator $E_n$.
Restriction and Yoneda give
$P_{n+1}/\m^{n+1}P_{n+1}\simeq P_n$.
As in C03, $P=\varprojlim P_n$ is finite over $\widehat A$ and represents
the restriction of $F$ on finite $\widehat A$-modules by the assumed
completion comparison, then on all such modules by filtered colimits.
C02 descends a finite representative to $A$. Necessity is exactly
Hom left exactness, finite presentation, completeness, and
$\Hom_{\widehat A}(P_{\widehat A},M)\subset M^r$ for a finite generating
set of $P_{\widehat A}$.

\section{94.C05b: the preliminary Artinian prorepresentability step}
\label{C05b}
\textbf{His anchor:} scan 6, \emph{En vertu de (i)} and
\emph{prorepr\'esentable}. Let $R$ be Artinian local and $G$ a
left-exact functor on finite $R$-modules, with its canonical
$R$-linear structure. It is strictly prorepresented there by finite
$R$-modules, without assuming that $G(E)$ is finite. If $G$ is defined
on all modules and preserves filtered colimits, this representation
holds on all modules.

\textbf{Proof (ours).} Put $D=\Hom_R(-,E)$ and $N=G(E)$, possibly
infinite. For finite $M$, the map
$G(M)\to\Hom_R(DM,N)$ sends $x$ to
$(\lambda\mapsto G(\lambda)x)$. It is an isomorphism for every
$E^r$ by $\End(E)=R$. A copresentation $0\to M\to E^r\to E^s$
then shows it is an isomorphism for every finite $M$, by comparing
kernels. Write $N$ as the directed union of its finite submodules
$N_j$. The module $DM$ is finite over the Artinian, hence noetherian,
ring $R$, and is finitely presented. Therefore
\[
 G(M)=\Hom_R(DM,N)=\colim_j\Hom_R(DM,N_j)
       =\colim_j\Hom_R(DN_j,M).
\]
The last equality is finite Artinian duality. Inclusions $N_j\subset
N_k$ dualize to surjections $DN_k\twoheadrightarrow DN_j$; thus the
prorepresentation is strict and all its terms are finite. If filtered
colimits are preserved, both functors extend this equality from finite
to arbitrary modules. When $N$ is finite the single representative
$DN$ gives C05. This verifies the manuscript's preliminary assertion,
not only the stronger conclusion obtained after the finiteness test.

\section{94.RD01: the smallest finiteness test in Proposition 4}
\label{RD01}
\textbf{Origin:} the parenthesis on scan 5 and the changed finite-length
condition on scan 6. This strengthening is ours, not attributed to him.
Under left exactness, filtered-colimit preservation and the completion
comparison of C05, finite representability is equivalent simply to
$\dim_k F(k)<\infty$.

\textbf{Proof (ours).}
For $0\to L'\to L\to k\to0$, left exactness gives
$0\to F(L')\to F(L)\to F(k)$, so finite length of the outside terms
implies finite length of $F(L)$, with
$\ell(F(L))\leq\ell(F(L'))+\ell(F(k))$.
A composition series therefore gives
$\ell(F(L))\leq\ell(L)\dim_k F(k)$ for every finite-length $L$.
In particular $F(E_n)$ is finite for every Artinian quotient's
injective cogenerator. The Artinian step and completion argument of
C05 now apply without further finiteness tests. Necessity follows
from $F(k)=\Hom_A(P,k)$ for finite $P$.
No open extension is asserted: the direction is completely settled.

\section{94.C06: Lemme 5 --- certified background}
\label{C06}
\textbf{His anchor:} scans 6--7, \emph{Alors}, replacing the deleted
claim of equivalence. For a noetherian ring $A$, assume left exactness,
filtered-colimit preservation, the completion comparison at every
$A_\p$ for finite $\widehat{A_\p}$-modules, finiteness on finite-length
$A_\p$-modules, and generic representability on every integral quotient.
The last condition must mean: for every prime $\p$, there is a
\emph{nonzero} $f\in A/\p$ such that $F_{(A/\p)_f}$ is representable.
Then $F$ is strictly prorepresented by finite modules $(P_i)$, and
for every $\p$ the localized pro-object is stationary and represents
$F_\p$. Here strict means that the transition maps of the representing
modules can be chosen surjective. Stationarity is an assertion about
the pro-object, not an assertion that one uniform index works on all
of $\Spec A$.

SETTLED H035.8 certifies this corrected sufficient criterion and its
minimal-support argument; it is carried, not independently reworked.
The manuscript does not write $f\ne0$ in (v). Without it, (v) is empty
of content, since one can take $f=0$ and the zero ring. This is his
omission, not permission to insert words into the transcription.
The certified proof uses the finite local representatives from C05,
their image submodules for a given element of $F(M)$, spreading at a
generic point of the remaining closed set, and Nakayama to extend the
open where this image is coherent. The boxed and left-margin additions
on scan 7 provide the intermediate containments and the reduction
$N'=M$. These are accounted for in \texttt{SCAN\_CHECKS.tsv}.

\section{94.C06b: why the missing nonzero qualification is essential}
\label{C06b}
\textbf{Origin:} Lemme 5(v), scan 6. The correction $f\ne0$ is certified
in H035.8. Here is our additional proof that dropping it actually
falsifies the stated sufficiency, rather than merely damaging the
written argument. Use the certified example of C07:
$A=\mathbb Z$, $L=\sum_{p\text{ prime}}p^{-1}\mathbb Z\subset\mathbb Q$,
and $F=L\otimes_{\mathbb Z}-$; all of (i)--(iv) hold by H035.10.
With $f=0$, condition (v) says only that the restriction to modules
over the zero ring is representable. That category has one object,
and $F(0)=0$, so it is representable by zero.

A strict prorepresentation by finite modules would give
$F(M)=\colim_i\Hom(P_i,M)$ with surjective transition maps
$P_j\twoheadrightarrow P_i$. Any element $\alpha$ represented by
$u:P_i\to M$ would then have a smallest supporting submodule,
namely $\im u$. Indeed it is supported there. If it is also supported
on $N\subset M$, represent that element by $v:P_j\to N$; at a common
later index the two maps into $M$ agree. Surjectivity onto $P_i$
forces $\im u\subset N$. This proves minimality and also its
independence of the chosen representation.

Apply this to $M=\mathbb Z$ and $\alpha=1\in L=F(\mathbb Z)$.
For every prime $p$, the image of $F(p\mathbb Z)$ in $F(\mathbb Z)$
is $pL$ and contains $1$, since $1/p\in L$. A smallest supporting
submodule would therefore be contained in
$\bigcap_{p\text{ prime}}p\mathbb Z=0$; but $1$ is not in the image
of $F(0)=0$. This contradiction rules out even a strict finite
prorepresentation, before imposing local stationarity. Thus the
literal theorem is false. The nonzero repair is the corrected theorem
C06, whose sufficiency is carried from H035.8 as instructed.

\section{94.C07: Remarque 5 bis --- certified refutation and repair}
\label{C07}
\textbf{His anchor:} scan 8, \emph{$F(N)$ est de type fini}.
Under C06(i)--(iv), and assuming that every integral quotient $B$ has
a nonempty open normal locus, he asserts equivalence between C06(v)
and constructibility of
$\p\mapsto\dim_{k(\p)}F(k(\p))$.
SETTLED H035.10 certifies the forward implication, refutes the
converse, and certifies its repair. The following is the supplied
counterexample, recorded as background, not claimed as a new proof:
\[
 A=\mathbb Z,\qquad L=\sum_{p\text{ prime}}\tfrac1p\mathbb Z
 \subset\mathbb Q,\qquad F(M)=L\otimes_{\mathbb Z}M.
\]
$L$ is torsion-free and hence flat; at a prime $(p)$ its localization
is $p^{-1}\mathbb Z_{(p)}$, free of rank one. At the generic point it
is $\mathbb Q$. Thus the local restrictions of $F$ are all represented
by a free rank-one module. In particular all local completion and
finite-length conditions hold, and its residue-field rank is constantly
one. Every integral quotient of $\mathbb Z$ has the required normal
open. Yet no $L[1/n]$ is finite over $\mathbb Z[1/n]$: finitely many
generators involve only finitely many denominator primes, whereas
$1/q$ occurs for every prime $q\nmid n$. If its tensor functor were
represented by a finite module on that open, evaluation on the ring
would make $L[1/n]$ finite. The missing condition is consequently
substantive.

\textbf{Certified correction.} Impose also finiteness of $F(B)$ for
every integral quotient $B$ (it is enough on a dense localization of
each such $B$). Then the converse holds; the normal-locus assumption
is unnecessary for this repair. The argument recorded in SETTLED is
that on a constant-rank open in an integral quotient, each finite local
representative $Q_\p$ has its number of generators equal to its generic
rank. A surjection from a free module of that rank has torsion kernel
inside a free module over a domain, hence zero kernel. Thus $Q_\p$ is
free. The natural map $F(B)\otimes_B-\to F(-)$ is an isomorphism on
these localizations. The added finiteness makes this a finite, locally
free tensor functor on the open, represented by the dual of $F(B)$.
This is the certified repair, not an alteration of his written
hypotheses. His simultaneous lattice inclusions are informal local
notation; our formulation avoids identifying an $A_\p$-module with
a finite global lattice.

\section{94.C08: Lemme 6, generic regular-sequence acyclicity}
\label{C08}
\textbf{His anchor:} scan 9, \emph{$\operatorname{gr}_J(F)$ loc. libres}.
\textbf{Modern statement (ours):} if $S$ is locally noetherian and
$\mathcal F$ coherent, there is a dense open $U\subset S$ on which
both $\operatorname{gr}_J\mathcal O_S$ and
$\operatorname{gr}_J\mathcal F$ are locally free of finite rank over
$\mathcal O_{S_{\rm red}}$, where $J$ is the nilradical. For every
$s\in U$ and every $A$-regular sequence $f_1,\ldots,f_n$ in the
maximal ideal of $A=\mathcal O_{S,s}$,
\[
 \Tor_i^A(\mathcal F_s,A/(f_1,\ldots,f_n))=0\quad(i>0).
\]
The local freeness concerns the finite direct sum of graded pieces;
the nilradical is nilpotent locally. No reducedness of $S$ is assumed.

\textbf{Existence of $U$ (ours).}
For a reduced noetherian affine scheme and finitely many finite
modules, there is a dense open on which they are locally free.
Indeed, near each of the finitely many generic points remove the
other components; this leaves a domain. Choose a vector-space basis
of each module over its fraction field. Lifting it to the module
and inverting finitely many nonzero elements makes both the kernel
and cokernel of the resulting free-module map zero. The union of
these opens contains all generic points. Apply this to the finitely
many graded pieces of $J$ and $\mathcal F$. For a locally noetherian
scheme, the locus where both graded objects are locally free is open
and is dense on every noetherian affine chart by this argument; hence
it is dense in $S$.

\textbf{Filtered algebra lemma and proof (ours).}
Let $A$ be noetherian local, $J$ a nilpotent ideal, $A_0=A/J$.
Suppose every graded piece of $A$ and a finite module $M$ is free over
$A_0$. If $f$ is a nonzerodivisor on $A$, its image $f_0$ is a
nonzerodivisor on $A_0$: choose the largest $r$ with $J^r\ne0$.
Then $J^r=\operatorname{gr}_J^r A$ is a nonzero free $A_0$-module
inside $A$, and multiplication by $f$ on it is multiplication by
$f_0$. Injectivity forces injectivity on $A_0$.
It follows that $f$ is injective on every finite filtration quotient
$M/J^qM$ as well as on $M$: if a nonzero class were killed, its first
nonzero graded class would be killed by $f_0$. Consequently
$fM\cap J^qM=fJ^qM$ for every $q$, and
\[
 \operatorname{gr}_{J'}(M/fM)
 =\operatorname{gr}_J M/f_0\operatorname{gr}_J M,
 \qquad J'=(J+fA)/fA.
\]
The same holds for $A/fA$. These graded modules are free over
$A_0/f_0A_0$. Repeating the argument proves that any $A$-regular
sequence is $M$-regular. It also identifies the graded module of the
successive quotient as a free module over
$A_0/(f_{10},\ldots,f_{n0})$. This induction uses only nilpotence of
$J'$; it does not require $A_0/(f_{10},\ldots,f_{n0})$ to remain reduced.

Finally the Koszul complex of an $A$-regular sequence is a free
resolution of its quotient: for one element this is the injective
two-term multiplication complex; inductively its tensor product with
the next two-term complex is the mapping cone of multiplication on
the previous resolution, whose sole homology is killed with zero
kernel by regularity. Tensoring with $M$ gives the same acyclicity
because the sequence is $M$-regular. This proves the Tor assertion in
all positive degrees, not just degree one.

\section{94.C08b: the shorthand in the proof of Lemme 6}
\label{C08b}
The live $M$ on scan 9 is not defined in the live statement; the
displayed definition of a different $M$ is struck out. Our proof
explicitly sets $M=\mathcal F_s$. His images $f_{i0}$ act on graded
modules, while the lifts $f_i$ act on $M$; these actions must not be
identified. Likewise the final assertion that
$\operatorname{gr}_J(N)$ is free needs the quotient coefficient ring
$A_0/(f_{10},\ldots,f_{n0})$. Over $A_0$ itself it is false: take
$A=k[t]_{(t)}$, $J=0$, $M=A$, $f=t$. Then $N=k$ is not free over
$A_0=A$ since it is nonzero and annihilated by $t$.
The filtered lemma in C08 proves the corrected statement completely.
These are explicit reconstructions of shorthand, not scan changes.

\section{94.C09: Corollaire 7 --- certified background}
\label{C09}
\textbf{His anchor:} scans 9--10, \emph{l'hom. naturel}.
For coherent $P,\Omega$ on a noetherian scheme $S$, there exists a
dense open $U$ such that for every $s\in U$ and every regular sequence
$\mathbf f$ in $A=\mathcal O_{S,s}$, the natural map
\[
 \Hom_A(P_s,\Omega_s)\otimes_A A/(\mathbf f)
 \longrightarrow\Hom_A(P_s,\Omega_s/(\mathbf f)\Omega_s)
\]
is an isomorphism. This conclusion is certified by SETTLED H035.15
and its Corollaire 7 note verdict. The two errors in its written proof
are separately recorded in C10 and C11; they do not refute the
corollary. With the corrected Hom complex, apply C08 to its finite
boundary and first-cohomology modules, and use the corrected
base-change criterion below. This records the certified dependency
without redoing the certified result.

\section{94.C10: the Hom-complex slip --- certified correction}
\label{C10}
\textbf{His anchor:} scan 10, the two factors \emph{$\otimes P$}.
SETTLED H035.14 confirms the ink and the error. For a finite free
presentation $L_1\to L_0\to P\to0$, the correct complex is
\[
 [L_0^\vee\otimes Q\longrightarrow L_1^\vee\otimes Q]
\]
in degrees zero and one; its zeroth cohomology is $\Hom(P,Q)$.
This follows by applying the left-exact Hom functor to the presentation
and identifying $\Hom(L_i,Q)=L_i^\vee\otimes Q$ for finite free $L_i$.
The supplied counterexample to his written factors is $A=k$, $P=k$,
$Q=k^2$, $L_0=k$, $L_1=0$: his complex has $H^0=k$, whereas
$\Hom(P,Q)=k^2$. Witness H silently makes the mathematical repair;
\texttt{SCAN\_CHECKS.tsv} restores the written $P$ in H. The corrected formula is
stated here, not smuggled into the reading.

\section{94.C11: the $Z^1$ criterion --- certified refutation}
\label{C11}
\textbf{His anchor:} scan 10, \emph{$Z^1(K)$}.
For a two-term complex $K=[K^0\to K^1]$, the two vanishings
$\Tor_1(M,H^1K)=\Tor_1(M,Z^1K)=0$ do \emph{not} imply that
$H^0(K)\otimes M\to H^0(K\otimes M)$ is an isomorphism.
SETTLED H035.15 supplies and certifies the counterexample
\[
 A=k[x,y]/(xy),\quad M=A/(x),\quad K=[A\xrightarrow{y}A].
\]
Here $Z^1=A$, $H^1=A/(y)$ and
$\Tor_1(A/(x),A/(y))=0$, but $H^0(K)=(x)$ and
$(x)\otimes A/(x)\simeq k$, whereas multiplication by $y$ is
injective on $A/(x)=k[y]$, so the target is zero. The Tor vanishing
can be read from the supplied periodic resolution with alternating
maps $x,y$: after tensoring with $A/(y)=k[x]$, the relevant kernel
of multiplication by $x$ is zero. All these checks are carried as
certified background. Replacing $Z^1$ by $B^1=\im(K^0\to K^1)$
gives the certified sufficient criterion. The next section supplies
an exact, necessary-and-sufficient refinement of that correction.

\section{94.RD02: exact obstructions to this base change}
\label{RD02}
\textbf{Origin:} the false sufficient criterion in C11. The refinement
here is ours. Let $K=[U\xrightarrow d V]$ be any two-term complex of
$A$-modules, let $H=\ker d$, $I=\im d$, $C=\coker d$, and let $M$
be any $A$-module. No flatness or finiteness assumption is required.
There is a natural exact sequence
\[
 0\to T_0\to H\otimes M\longrightarrow H^0(K\otimes M)
 \to T_1\to0,
\]
where
\[
 T_0=\coker[\Tor_1(M,U)\to\Tor_1(M,I)],\qquad
 T_1=\coker[\Tor_1(M,V)\to\Tor_1(M,C)].
\]
Thus the comparison is an isomorphism \emph{exactly} when both
indicated Tor maps are surjective. Vanishing of
$\Tor_1(M,I)$ and $\Tor_1(M,C)$ is sufficient, and is the corrected
version closest to his statement. This exact criterion is stronger
than merely substituting one letter.

\textbf{Proof (ours).}
Tensor $0\to H\to U\to I\to0$. Its long exact sequence identifies
the kernel of $H\otimes M\to U\otimes M$ with $T_0$ and its image
with the kernel of the surjection $U\otimes M\to I\otimes M$.
Tensor $0\to I\to V\to C\to0$. Its long exact sequence identifies
the kernel of $I\otimes M\to V\otimes M$ with $T_1$.
The kernel of the composite $U\otimes M\to V\otimes M$ maps
surjectively onto this last kernel because the first map is
surjective. Its own kernel is exactly the preceding image of
$H\otimes M$. These three observations give the displayed sequence,
including the endpoint zeroes and its naturality.
If $U,V$ are flat, this specializes to
\[
 0\to\Tor_2(M,C)\to H\otimes M\to H^0(K\otimes M)
 \to\Tor_1(M,C)\to0,
\]
by dimension shifting through $0\to I\to V\to C\to0$.
In the counterexample of C11 the failure is precisely the nonzero
left obstruction; checking only the right obstruction misses it.
The direction is settled; no open extension is asserted.

\section{94.C12: the preliminary equivalence on scan 11}
\label{C12}
\textbf{His anchor:} \emph{module de diff\'erentielles pour un C.M.}.
The ink really says this; it is not silently replaced by
\emph{dualizing module}. No coefficient field or differential degree
is given. The neighboring Lemme 8 and Corollaire 9 concern a
\emph{faisceau fondamental}. Our precise interpretation of the
equivalence is therefore the following local canonical-module
statement, with its missing quantifiers explicitly supplied by us.

Let $(A,\m,k)$ be a noetherian Cohen--Macaulay local ring and $W$ a
finite $A$-module. Then $W$ is a dualizing module if and only if for
every system of parameters $\mathbf x$, the module
$W/(\mathbf x)W$ is the injective hull of $k$ over $A/(\mathbf x)$.
It is enough to require this for the powers
$(x_1^n,\ldots,x_d^n)$, $n\ge1$, of a single fixed system of parameters.
Here dualizing module means that $W[0]$ is a dualizing complex;
the normalized complex is then $W[d]$, $d=\dim A$.

\textbf{Proof of this reconstruction (ours).}
The forward implication is C13--C14 below. For the converse put
$\mathfrak q_n=(x_1^n,\ldots,x_d^n)$ and $R=\widehat A$.
The complete local ring $R$ has a dualizing complex
(Stacks Lemma 47.22.4, \cite{StacksDualizingComplete}), and is CM,
so it has a dualizing module $V$ (Lemma 47.20.2,
\cite{StacksDualizingCM}). By the forward implication,
$V/\mathfrak q_nV$ and $W/\mathfrak q_nW$ are both injective hulls
over $A/\mathfrak q_n$. Choose an isomorphism between them for $n=1$.
Given one for $n$, an arbitrary isomorphism at level $n+1$ induces
another at level $n$. The ratio is an automorphism of the injective
hull at level $n$. Its endomorphism ring is $A/\mathfrak q_n$
(Artinian duality, \cite{StacksArtinDuality}), so this ratio is
multiplication by a unit. That unit lifts to a unit modulo
$\mathfrak q_{n+1}$, since these are local rings. Multiplying the
chosen next isomorphism by that lift makes it compatible. Induction
produces a compatible system.

The ideals $\mathfrak q_n$ are cofinal with the $\m$-adic powers:
they are contained in $(\mathbf x)^n$, and
$(\mathbf x)^{d(n-1)+1}\subset\mathfrak q_n$, while $(\mathbf x)$
is $\m$-primary. Taking inverse limits therefore gives
$\widehat W\simeq V$. Faithfully flat descent of a dualizing complex
(Stacks Lemma 47.26.1, Tag 0E4A, PDF pp. 51--52,
\cite{StacksDualizingDescent}) applied to $W[0]$ shows that $W$ is
dualizing over $A$. For $d=0$ the empty parameter system gives the
assertion directly by Artinian duality; no power argument is needed.
This proof supplies compatible choices rather than assuming that
unrelated quotient isomorphisms automatically pass to the limit.

\textbf{Literal K\"ahler-differential interpretation (refuted).}
If the heading were instead read as a claim about $\Omega^1_{A/k}$,
it would be false. For $A=k[t,u]_{(t,u)}$, a regular CM local ring,
$\Omega^1_{A/k}$ is free of rank two. Modulo the parameter ideal
$(t,u)$ it is $k^2$, whereas the indecomposable injective over the
residue field is $k$. The canonical module is rank one instead.
This counterexample decides the literal modern interpretation;
it is not evidence for changing the historical wording. The
canonical-module reconstruction above is fully proved.

\section{94.C13: Lemme 8, quotient of a fundamental sheaf}
\label{C13}
\textbf{His anchor:} scan 11, \emph{faisceau fondamental}.
Our modern definition of that term here is a coherent sheaf whose
stalk at every point is a dualizing module over the CM local ring.
Let $X$ be locally noetherian and CM, with such a sheaf $\Omega$,
and let $\mathbf f=(f_1,\ldots,f_r)$ be an
$\mathcal O_X$-regular sequence of sections. Then
$Y=V(\mathbf f)$ is CM and $\Omega/(\mathbf f)\Omega$ is a
fundamental sheaf on $Y$. The assertion is stalkwise and is vacuous
off $Y$.

\textbf{Literature match and calculation.}
At a stalk $(A,\m)$, the normalized dualizing complex is
$\omega_A[d]$ with $\omega_A$ maximal CM by Stacks Lemma 47.20.2
(Tag 0AWS; dualizing chapter p. 36). A regular sequence in the
CM ring is a partial parameter system and is regular on this
maximal CM module. The Koszul calculation therefore gives
\[
 R\Hom_A(A/(\mathbf f),\omega_A[d])
 \simeq (\omega_A/(\mathbf f)\omega_A)[d-r].
\]
By Stacks Lemma 47.16.2 (Tag 0A7N) this is the normalized dualizing
complex of $A/(\mathbf f)$; the single-element form is Lemma
47.16.10 (Tag 0A7T, chapter p. 32). All these sources were opened
\cite{StacksDualizingLocal,StacksDualizingCM}.
Iteration proves the stated quotient property. The ordered global
sequence supplies the usual Koszul orientation; no global dualizing
complex or unmentioned choice of a base field is assumed.

\section{94.C14: Corollaire 9, the Artinian quotient}
\label{C14}
\textbf{His anchor:} scan 11, \emph{l'injectif ind\'ecomposable}.
For $x\in X$ as in C13 and a full system of parameters $\mathbf f$
of $A=\mathcal O_{X,x}$, the quotient
$\Omega_x/(\mathbf f)\Omega_x$ is the injective hull of $k(x)$ over
the Artinian local ring $A/(\mathbf f)$.
In a CM local ring a full parameter system is regular; C13 makes
this module dualizing. Stacks Lemma 47.16.8 (Tag 0A7R, chapter p. 32)
identifies the normalized dualizing complex of a zero-dimensional
local ring with that injective hull \cite{StacksDualizingLocal}.
This verifies the exact claim and all its hypotheses. The stalk
subscript omitted in his quotient notation is supplied explicitly.

\section{94.C15: Corollaire 10(a),(b), existence on a dense open}
\label{C15}
\textbf{His anchor:} scans 11 and 13, \emph{localemt immergeable dans un r\'egulier}.
Let $X$ be a locally noetherian scheme, locally isomorphic to a locally
closed subscheme of a regular locally noetherian scheme. There exist a
dense open $U\subset X$ and a coherent $\mathcal O_X$-module $\Omega$
such that $U$ is CM and, for every $x\in U$ and every full parameter
system $\mathbf f$ of $A=\mathcal O_{X,x}$,
$\Omega_x/(\mathbf f)\Omega_x$ is the injective hull of $k(x)$ over
$A/(\mathbf f)$. In particular $\Omega$ is on all of $X$, as his
statement requires, not just on $U$. The proof below is ours; it does
not assume that $X$ is quasi-compact or that all of $X$ has a global
dualizing complex.

\textbf{Proof.} First work near the generic point $\eta$ of one
irreducible component, after removing the other components. A locally
closed immersion can be made closed by shrinking its ambient scheme,
so choose an affine neighborhood $\operatorname{Spec}B$ of $\eta$,
$B=R/I$, with $R$ noetherian regular. Let $\mathfrak p\subset R$ be
the inverse image of $\eta$. The local ring $B_{\mathfrak p}$ is
Artinian. The regular local ring $R_{\mathfrak p}$ has dimension $c$
and global dimension at most $c$ (Stacks Proposition 10.110.1,
\cite{StacksRegularResolution}).
It is Gorenstein (Lemma 47.21.3, \cite{StacksRegularGorenstein}), and local duality for the finite
quotient gives
\[
 \Ext^i_{R_{\mathfrak p}}(B_{\mathfrak p},R_{\mathfrak p})=0
 \quad(i\ne c),
\]
while the degree-$c$ module is a dualizing module of $B_{\mathfrak p}$.
These are the exact local statements of Lemmas 47.16.2 and 47.20.2
already used in C13--C14.

A finite free resolution at $\mathfrak p$ spreads to a neighborhood:
take the first $c$ steps of a finite free resolution over $R$; its
last finite syzygy is free at $\mathfrak p$, and a basis and its inverse
spread after inverting one element outside $\mathfrak p$. Thus $B$
has a finite free resolution there of length $c$. Its finitely many
Ext modules in degrees different from $c$ vanish at $\mathfrak p$;
shrinking again kills all of them. Put $W=\Ext^c_R(B,R)$ on this
neighborhood. At any point of its intersection with $\operatorname{Spec}B$,
the local ambient ring is regular, hence dualizing over itself;
$R\Hom(B,R)$ is dualizing for the finite quotient and is concentrated
in degree $c$. Lemma 47.20.2 consequently makes that quotient ring
CM and $W$ its dualizing module. This argument uses no bound on the
dimensions of ambient points away from the chosen neighborhood.

The irreducible components of a locally noetherian scheme form a
locally finite family: a noetherian affine open meets only the finitely
many components whose generic points it contains. Thus, for each
component $Z_i$, the complement of the union of all $Z_j$, $j\ne i$,
is open and contains its generic point. Choose the preceding
neighborhood $V_i$ inside this complement. The $V_i$ are pairwise
disjoint. Their union $U$ is dense, and the modules $W_i$ therefore
define a coherent sheaf $W$ on $U$, with no gluing obstruction.

For completeness, $W$ extends coherently to $X$. On an affine
noetherian scheme $T$ a coherent sheaf $G$ on an open $T'$ has such
an extension: $T'$ is quasi-compact, so $j_*G=\widetilde M$ is
quasi-coherent (Stacks Lemma 28.23.1, \cite{StacksSheafExtension}). Cover $T'$ by finitely many
principal opens and choose finitely many numerators in $M$ that
generate the localizations on each; their finite submodule $N$
satisfies $\widetilde N|_{T'}=G$. Noetherianity makes $\widetilde N$
coherent. Well-order an affine open covering of $X$. Successively
extend the sheaf from the intersection with the next affine by this
construction, and glue; at a limit stage take the compatible union
on the union of the preceding opens. Coherence is local, so every
stage remains coherent. This gives a coherent extension $\Omega$
on $X$. Finally C14 supplies all parameter-quotient properties on
$U$. The empty scheme case is immediate. This proves the full stated
scope, including the global extension of the sheaf.

\section{94.C16: the marginal weakening of Corollaire 10 --- certified verdict}
\label{C16}
\textbf{His anchor:} scan 11 margin, \emph{ou seulement} and
\emph{EGA IV 6.11.8}. The supplied SETTLED H035.21 certifies that
excellent schemes remain allowed, but the further CMU weakening
is false. CMU here means that every closed integral subscheme has
a nonempty CM open; it does not supply a coherent canonical module
on a dense open. The certified counterexample is the localization
of
\[
 k[t_i^3,t_i^4,t_i^5\mid i\ge1]
\]
away from the union of the block ideals. The settled record
checks that it is a one-dimensional noetherian CM domain and that
non-Gorenstein points meet every nonempty open. A finite canonical
module would have generic rank one and become free on a smaller
nonempty open, forcing that open to be Gorenstein, a contradiction.
The construction's noetherianity and the semigroup calculation are
\emph{carried as certified background}; they have not been recomputed
or promoted to an independent proof here. The closest corrected
margin keeps the excellent clause and removes the asserted CMU
sufficiency. C15 supplies a complete independent proof under the
main text's original immersion hypothesis. This does not identify
excellence with that hypothesis or infer one from the other.

\section{94.C17: the definition of absolute value --- certified correction}
\label{C17}
\textbf{His anchor:} scan 12, \emph{$f(x)=0$ entra\^ine $x=0$}.
The printed definition requires multiplicativity, the ordinary triangle
inequality, and only that implication. SETTLED H035.22 certifies that
$f\equiv1$ satisfies every written axiom, but $f(x-x)=1$, so the
asserted distance is not a distance. The closest corrected definition
adds $f(0)=0$; then $f(x)=0$ exactly at zero. With that correction
$d(x,y)=f(x-y)$ is a metric and gives a field topology. This corrected
result, including continuity of addition, multiplication and inversion,
is carried as the certified background H035.22, not independently
re-proved here. All later claims on this leaf are interpreted with that
normalization. Its absence is an authorial slip, not a scan error.

\section{94.C18: exponentiating a real valuation}
\label{C18}
\textbf{His anchor:} scan 12, \emph{$e^{-v(x)}$}.
His word \emph{archim\'edien} is immediately defined to mean an ordered
value group isomorphic to an additive ordered subgroup of $\mathbb R$.
Choose such an order-preserving embedding. Our explicit conventions
are $v(0)=+\infty$ and $e^{-\infty}=0$; on $K^\times$, $v$ is
additive and $v(x+y)\ge\min(v(x),v(y))$ whenever the sum is nonzero.
Then $f(x)=e^{-v(x)}$ is a normalized absolute value, satisfying the
strong triangle inequality. Conversely every normalized absolute
value with that inequality arises this way.

\textbf{Proof (ours).} For nonzero $x,y$, exponentiation changes
$v(xy)=v(x)+v(y)$ into multiplicativity. The valuation inequality
changes into $f(x+y)\le\max(f(x),f(y))\le f(x)+f(y)$; the cases
involving zero follow from the conventions. Nonzero values are
strictly positive. Conversely set $v(x)=-\log f(x)$ on $K^\times$.
Multiplicativity gives additivity, and the strong triangle inequality
gives the valuation inequality after applying the decreasing function
$-\log$. Its value group is the image subgroup of $\mathbb R$ with
the induced order. The trivial absolute value corresponds to the
zero group. No stronger embedding theorem is being assumed: the
embedding is part of his explicit definition on this sheet.

\section{94.C19: Proposition 1, the binomial criterion}
\label{C19}
\textbf{His anchor:} scan 12, \emph{l'une ou l'autre des conditions}.
For a \emph{normalized} absolute value $f$ on a field, the following
are equivalent: it arises from the real valuation in C18; it satisfies
$f(x+y)\le\max(f(x),f(y))$; and $f(n\cdot1)\le1$ for every natural
integer $n$. Without the normalization from C17 the literal proposition
also fails for $f\equiv1$, which satisfies (a) and (b) but cannot be
$e^{-v}$ at zero.

\textbf{Proof of the corrected statement (ours).} C18 proves the
first equivalence. Multiplicativity and positivity give $f(1)=1$.
Repeated use of the strong triangle inequality gives $f(n\cdot1)\le1$;
$n=0$ uses normalization. Conversely, let $R=\max(f(a),f(b))$.
If $R=0$ then $a=b=0$ and the desired inequality holds. For $R>0$
and any integer $n\ge1$, expand $(a+b)^n$. Each binomial coefficient
is a natural integer acting through the prime subring, so its value
is at most one. Each of the $n+1$ summands has absolute value at most
$R^n$. The ordinary triangle inequality therefore gives
\[
 f(a+b)^n=f((a+b)^n)\le(n+1)R^n.
\]
Taking positive $n$th roots and using
$\log(n+1)/n\to0$ gives $f(a+b)\le R$. This checks both the
printed coefficient $n+1$ and its limiting use, in every characteristic.

\section{94.C20: the positive-characteristic corollary --- certified correction}
\label{C20}
\textbf{His anchor:} scan 12, \emph{tout \'el\'ement $n.1$}.
The corollary for normalized absolute values is true. Its stated reason
that every $n\cdot1$ is a root of unity is false at $n=p$ in
characteristic $p$. SETTLED H035.24 certifies the correction: every
\emph{nonzero} element of the prime field has finite multiplicative
order; hence its absolute value is one, and zero has value zero.
Condition (b) of C19 follows. The statement and counterexample are
carried from that settled result; the ink was separately checked.

\section{94.C21: equivalence of absolute values --- certified correction}
\label{C21}
\textbf{His anchor:} scan 12, \emph{un nombre r\'eel $s$}.
SETTLED H035.25 certifies the missing restriction: equivalence of
normalized absolute values is $f=g^s$ with $s>0$, and is equivalent
to equality of their field topologies. At zero a negative exponent
is not defined as a finite value; exponent zero either leaves $0^0$
undefined or, if values on $K^\times$ are extended by zero, turns any
nontrivial absolute value into the trivial one. For example on
$\mathbb Q$ the ordinary absolute value has $1/n\to0$, whereas the
trivial value gives the discrete topology. Thus zero is not an allowed
equivalence exponent. The fully corrected topological equivalence
is certified background, not an independently claimed proof here.
The Bourbaki citation is transcribed as a historical reference; no
unopened edition is used to certify this result.

\section{94.C22: Proposition 2 on $\mathbb Q$ --- certified correction}
\label{C22}
\textbf{His anchor:} scan 12, \emph{$p$-adique, ou ... ordinaire}.
The trivial normalized absolute value is missing. It is neither a
positive power of a $p$-adic value (evaluate at $p$) nor a positive
power of the ordinary one (evaluate at $2$). SETTLED H035.26 certifies
the full corrected classification: a normalized absolute value on
$\mathbb Q$ is trivial, or $|\cdot|_p^s$ for a unique prime $p$ and
$s>0$, or $|\cdot|_\infty^s$ with $0<s\le1$. The restriction $s\le1$
in the last case is required by the ordinary triangle inequality,
as $2^s\le2$; the classification and its corrected proof are carried
from the supplied settled record, as instructed.

His unfinished proof on this leaf begins with integers $a,b>1$ and
the base-$b$ expansion $a^h=\sum_{i=0}^q a_i b^i$,
$0\le a_i\le b-1$. This is valid by Euclidean division. In the
non-valuative case C19 ensures an integer $n$ with $f(n)>1$.
The printed last words have no following estimate on this scan.
No continuation is supplied as his: scan 13 resumes Corollaire 10,
not this proof. The settled classification decides the mathematical
item without reconstructing an absent page or redoing its certified proof.

\section{94.C23: Corollaire 10(c), detection by parameter-quotient duals}
\label{C23}
\textbf{His anchor:} scan 13, \emph{pour tt syst. de param\`etres}.
Let $(A,\m,k)$ be noetherian CM local and let $\Omega$ be finite with
$E_{\mathfrak q}=\Omega/\mathfrak q\Omega$ the injective hull of $k$
over $A/\mathfrak q$ for every parameter ideal $\mathfrak q$.
If $u:R\twoheadrightarrow R'$ is a surjection of finite modules
such that
\[
 \Hom_A(R',E_{\mathfrak q})\longrightarrow
 \Hom_A(R,E_{\mathfrak q})
\]
is an isomorphism for every such ideal, then $u$ is an isomorphism.
This is his written claim with stalks made explicit. Our stronger
form drops the assumed surjectivity and tests only
$\mathfrak q_n=(x_1^n,\ldots,x_d^n)$ for one full parameter system
and all $n\ge1$.

\textbf{Proof (ours).} Every map to $E_{\mathfrak q_n}$ factors through
the corresponding quotient of its source. Artinian duality over
$A/\mathfrak q_n$ is a faithful exact duality on finite modules
(C14 and Stacks Lemmas 47.6.2--47.6.3), so the assumed dual comparison
is an isomorphism if and only if
$R/\mathfrak q_nR\to R'/\mathfrak q_nR'$ is an isomorphism.
For $n=1$, right exactness and Nakayama applied to the finite cokernel
show that $u$ is surjective even if that was not assumed. If
$z\in\ker u$, injectivity of every quotient comparison implies
$z\in\mathfrak q_nR$ for every $n$. These ideals are cofinal with
the powers of $\m$ as proved in C12, so Krull intersection for the
finite module $R$ gives $z=0$. One elementary form of that last step
is Artin--Rees: for $N=\bigcap_n\m^nR$, finite by noetherianity,
Artin--Rees implies $N=\m N$, and Nakayama gives $N=0$.
If $d=0$, the empty parameter ideal is zero and the single Artinian
duality test already proves the assertion. This also shows why testing
only one nonzero parameter ideal would be insufficient: over a DVR,
$A\to A/(t)$ induces an isomorphism modulo $(t)$ but is not injective.
The cofinal family, not a single test, is used by the proof.

\section{94.C24: the canceled Corollaire 11}
\label{C24}
\textbf{His anchor:} scan 13, \emph{Sous ces conditions}; the entire
draft is canceled. With $A$, $\Omega$ as in C23, let
$u:R\twoheadrightarrow R'$ be a map of finite $A$-modules. Define
$\Phi_R(M)=\Hom_A(R,M)$, as in his rewritten formulas, and similarly
for $R'$. Suppose, for every $A$-regular sequence $\mathbf f$ in the
maximal ideal and $T=R,R'$, that
\[
 \Hom_A(T,\Omega)\otimes A/(\mathbf f)
 \longrightarrow\Hom_A(T,\Omega/(\mathbf f)\Omega)
\]
is an isomorphism. Suppose also that the map induced by $u$,
$\Hom_A(R',\Omega)\to\Hom_A(R,\Omega)$, is an isomorphism.
Then $u$ is an isomorphism. Our weaker sufficient hypotheses require
the base-change comparisons only for the cofinal parameter powers
of C23 and do not require surjectivity in advance.

\textbf{Proof (ours).} Tensor the second assumed isomorphism with
$A/\mathfrak q_n$. Naturality of the two base-change comparisons
identifies this tensor map with
$\Hom_A(R',E_{\mathfrak q_n})\to\Hom_A(R,E_{\mathfrak q_n})$.
Both vertical maps and the tensor map are isomorphisms; therefore
the latter Hom map is an isomorphism. Parameter powers are regular
because $A$ is CM. Apply C23. This proof is independent of whether
the canceled earlier functors contained $D\otimes M$ or
$\Omega\otimes M$: only the clearly rewritten final $\Hom(T,M)$
formulas are used. Cancellation is not a mathematical counterexample,
and the marginal request for a different exposition is not treated
as a proof or an unresolved theorem.

\bibliographystyle{plain}
\nocite{EGAIV2,BourbakiHistorical}
\bibliography{references}
\end{document}
